% ============================================================================
%  back/vocabulary.tex — واژه‌نامهٔ جامعِ الفباییِ کلِ جزوه
%  واژه‌های درس ۱ و ۲ به‌صورتِ کامل درج شده است؛ الگوهای درس ۳ تا ۵
%  آمادهٔ تکمیل است.
% ============================================================================
\Jpart{واژه‌نامهٔ جامعِ الفبایی}
  {اَلْمُعْجَمُ الْجَامِعُ}
  {همهٔ واژه‌های مُعجَمِ دروس، به ترتیبِ الفبا؛ با روابطِ معنایی (مترادف،
   متضاد، جمع و مفرد) و شمارهٔ درس}

\begin{JGTable}

% ------------------------------------------------------------- حرفِ ا -------
\JGLetter{ا / أ / إ}
\SJW{أَجْرَى}{اجرا کرد}{مضارع: يُجْرِي ← لِيُجْرِيَ: تا اجرا کند}{٢}
\SJW{أَبْصَرَ}{نگاه کرد}{}{١}
\SJW{أَضْعَفَ}{ضعیف کرد}{}{٢}
\SJW{أَفَادَ}{سود رساند}{مضارع: يُفِيدُ}{٢}
\SJW{أَقْبَلَ عَلَى}{به... روی آورد}{}{٢}
\SJW{اِهْتَمَّ}{اهتمام ورزید}{مضارع: يَهْتَمُّ}{٢}
\SJW{اَلْإِعْمَار}{آباد کردن}{أَعْمَرَ ، يُعْمِرُ}{٢}
\SJW{اَلْأَنْحَاء}{سمت‌ها، سوها}{مفرد: اَلنَّحْو}{٢}
\SJW{اَلْأَنْفَاق}{تونل‌ها}{مفرد: اَلنَّفَق}{٢}
\SJW{اِنْطَوَى}{به هم پیچیده شد، نهان شد}{مضارع: يَنْطَوِي}{١}
\SJW{اِنْهَدَمَ}{ویران شد}{}{٢}
\SJW{اَلْأَهْل}{شایسته}{}{٢}

% ------------------------------------------------------------- حرفِ ب -------
\JGLetter{ب}
\SJW{اَلْبَدَل}{جانشین، بهرهٔ مادی}{جمع: اَلْأَبْدَال}{١}

% ------------------------------------------------------------- حرفِ ت -------
\JGLetter{ت}
\SJW{اَلتَّحْوِيل}{دگرگونی}{}{٢}
\SJW{اَلتَّطْوِير}{بهینه‌سازی}{}{٢}
\SJW{اَلتِّقْنِيَّة}{فناوری (تکنیک)}{}{٢}
\SJW{اَلتِّلَال}{تپه‌ها}{مفرد: اَلتَّلّ}{٢}
\SJW{تَمَّ}{انجام شد، کامل شد}{مضارع: يَتِمُّ}{٢}

% ------------------------------------------------------------- حرفِ ج -------
\JGLetter{ج}
\SJW{اَلْجِرْم}{پیکر، تن}{جمع: اَلْأَجْرَام}{١}
\SJW{جُلْبُ شَعِيرَةٍ}{پوستِ جو}{}{٢}

% ------------------------------------------------------------- حرفِ ح -------
\JGLetter{ح}
\SJW{اَلْحَدِيد}{آهن}{}{١}
\SJW{حَدَّدَ}{مشخص کرد}{}{٢}

% ------------------------------------------------------------- حرفِ خ -------
\JGLetter{خ}
\SJW{خَيْبَةُ الْأَمَلِ}{ناامیدی}{\textbf{≠} اَلرَّجَاء}{٢}

% ------------------------------------------------------------- حرفِ د -------
\JGLetter{د}
\SJW{اَلدَّاء}{بیماری}{\textbf{=} اَلْمَرَض ؛ \textbf{≠} اَلشِّفَاء ، اَلصِّحَّة}{١}
\SJW{اَلدَّؤُوب}{با پشتکار}{}{٢}

% ------------------------------------------------------------- حرفِ ز -------
\JGLetter{ز}
\SJW{زَعَمَ}{گمان کرد}{}{١}

% ------------------------------------------------------------- حرفِ س -------
\JGLetter{س}
\SJW{سَهَّلَ}{آسان کرد}{\textbf{≠} صَعَّبَ}{٢}
\SJW{سِوَى}{جُز، غیر از}{}{١}
\SJW{اَلسُّهُول}{دشت‌ها}{مفرد: اَلسَّهْل}{٢}
\SJW{اَلسُّوَيْد}{سوئد}{}{٢}

% ------------------------------------------------------------- حرفِ ش -------
\JGLetter{ش}
\SJW{اَلشَّعِير}{جو}{}{٢}
\SJW{اَلشَّقّ}{شکافتن}{شَقَّ ، يَشُقُّ}{٢}

% ------------------------------------------------------------- حرفِ ص -------
\JGLetter{ص}
\SJW{اَلصَّالِحَةُ لِلزِّرَاعَةِ}{قابلِ کشت}{}{٢}
\SJW{اَلصَّبِيّ}{کودک، پسر}{جمع: اَلصِّبْيَان}{٢}
\SJW{صَحَّحَ}{تصحیح کرد}{}{٢}

% ------------------------------------------------------------- حرفِ ط -------
\JGLetter{ط}
\SJW{اَلطُّنّ}{تُن}{جمع: اَلْأَطْنَان}{٢}
\SJW{اَلطِّين ، اَلطِّينَة}{گِل، سرشت}{}{١}

% ------------------------------------------------------------- حرفِ ع -------
\JGLetter{ع}
\SJW{اَلْعَصَب}{پی، عصب}{جمع: اَلْأَعْصَاب}{١}
\SJW{اَلْعَظْم}{استخوان}{جمع: اَلْعِظَام}{١}

% ------------------------------------------------------------- حرفِ ف -------
\JGLetter{ف}
\SJW{اَلْفَرَنْسِيَّة}{زبانِ فرانسوی}{}{٢}
\SJW{اَلْفِيزِيَاء}{فیزیک}{}{٢}

% ------------------------------------------------------------- حرفِ ق -------
\JGLetter{ق}
\SJW{قَنَاةُ بَنَمَا}{کانالِ پاناما}{}{٢}
\SJW{اَلْقَنَوَات}{کانال‌ها}{مفرد: اَلْقَنَاة}{٢}

% ------------------------------------------------------------- حرفِ ك -------
\JGLetter{ك}
\SJW{كَسَبَ}{به دست آورد}{}{٢}

% ------------------------------------------------------------- حرفِ ل -------
\JGLetter{ل}
\SJW{اَللَّحْم}{گوشت}{جمع: اَللُّحُوم}{١}

% ------------------------------------------------------------- حرفِ م -------
\JGLetter{م}
\SJW{اَلْمَجَال}{زمینه}{جمع: اَلْمَجَالَات}{٢}
\SJW{اَلْمَنَاجِم}{معادن}{مفرد: اَلْمَنْجَم}{٢}

% ------------------------------------------------------------- حرفِ ن -------
\JGLetter{ن}
\SJW{نَشَرَ}{پخش کرد}{}{٢}
\SJW{اَلنُّحَاس}{مس}{}{١}

\end{JGTable}

% --------------------------------------- بخشِ درس‌های ۳ تا ۵ (الگو) --------
\Jsec{اَلْمُعْجَمُ — درسِ ٣ تا ٥ \,\textbar\, محلِ تکمیل}
\begin{JGood}[راهنمای تکمیل]
واژه‌های معجَمِ درسِ سوم (\emph{ثَلَاثُ قِصَصٍ قَصِيرَةٍ})، چهارم
(\emph{نِظَامُ الطَّبِيعَةِ}) و پنجم (\emph{يَا إِلٰهِي}) را پس از مطالعهٔ
هر درس، به همین جدول بیفزایید و شمارهٔ درس را در ستونِ آخر بنویسید.
\end{JGood}

\begin{JMTable}[اَلْمُعْجَمُ \,\textbar\, جدولِ تکمیلیِ درسِ سوم]
\JW{...}{...}{}
\JW{...}{...}{}
\end{JMTable}
\begin{JMTable}[اَلْمُعْجَمُ \,\textbar\, جدولِ تکمیلیِ درسِ چهارم]
\JW{...}{...}{}
\JW{...}{...}{}
\end{JMTable}
\begin{JMTable}[اَلْمُعْجَمُ \,\textbar\, جدولِ تکمیلیِ درسِ پنجم]
\JW{...}{...}{}
\JW{...}{...}{}
\end{JMTable}
